\subsection*{\texorpdfstring{Positivity on $\mathbb{Z}$}{Positivity on Z}}

\begin{definitionbox}{Definition}
\begin{definition}[Positive Integer]
\label{def:positive-integer}
An integer $x\in\mathbb{Z}$ is \emph{positive} if $x=[a,b]$ for some
$a,b\in\mathbb{W}$ with $b<a$.
\end{definition}
\end{definitionbox}

\begin{remark*}[Standard quantified statement]
\[
\operatorname{Positive}_{\mathbb{Z}}(x)
\Longleftrightarrow
\exists a,b\in\mathbb{W},\ x=[a,b]\land b<a.
\]
\end{remark*}

\begin{remark*}[Predicate reading]
\[
\operatorname{Positive}(x,\mathbb{Z},<_{\mathbb{Z}},0_{\mathbb{Z}}).
\]
\end{remark*}

\begin{remark*}[Interpretation]
Positive integers are represented by formal differences whose first coordinate
is larger than the second.
\end{remark*}

\begin{dependencies}
\begin{itemize}
  \item \hyperref[def:integer-class-notation]{Integer Class Notation}
  \item \hyperref[def:whole-numbers]{Whole Numbers}
  \item \hyperref[def:order-on-w]{Order on $\mathbb{W}$}
\end{itemize}
\end{dependencies}

\begin{lemmabox}{Lemma}
\begin{lemma}[Positivity Respects Formal-Difference Equivalence]
\label{lem:positivity-respects-equivalence}
If $[a,b]=[a',b']$, then
\[
b<a \Longleftrightarrow b'<a'.
\]
\hyperref[prf:positivity-respects-equivalence]{\textit{Go to proof.}}
\end{lemma}
\end{lemmabox}

\begin{remark*}[Standard quantified statement]
\[
[a,b]=[a',b']
\Longrightarrow
\bigl(b<a \Longleftrightarrow b'<a'\bigr).
\]
\end{remark*}

\begin{remark*}[Predicate reading]
\[
\operatorname{RepresentativeIndependent}(\operatorname{Positive}_{\mathbb{Z}},
\sim_{\mathbb{Z}},\operatorname{Pre}\mathbb{Z}).
\]
\end{remark*}

\begin{remark*}[Interpretation]
Positivity is independent of the chosen representative of an integer.
\end{remark*}

\begin{dependencies}
\begin{itemize}
  \item \hyperref[thm:addition-on-w-cancellation]{Addition on $\mathbb{W}$ Satisfies Cancellation}
  \item \hyperref[thm:order-on-w-total]{Order on $\mathbb{W}$ Is Total}
\end{itemize}
\end{dependencies}

\begin{theorembox}{Theorem}
\begin{theorem}[Positivity on $\mathbb{Z}$ Is Well-Defined]
\label{thm:positivity-on-z-well-defined}
The predicate ``$x$ is positive'' is well-defined on equivalence classes in
$\mathbb{Z}$.
\hyperref[prf:positivity-on-z-well-defined]{\textit{Go to proof.}}
\end{theorem}
\end{theorembox}

\begin{remark*}[Standard quantified statement]
\[
\operatorname{WellDefinedPredicate}(\operatorname{Positive}_{\mathbb{Z}},\mathbb{Z}).
\]
\end{remark*}

\begin{remark*}[Predicate reading]
\[
\operatorname{WellDefinedPredicate}(\operatorname{Positive}_{\mathbb{Z}},\mathbb{Z}).
\]
\end{remark*}

\begin{remark*}[Interpretation]
The positivity test descends from representatives to the integer quotient.
\end{remark*}

\begin{dependencies}
\begin{itemize}
  \item \hyperref[def:positive-integer]{Positive Integer}
  \item \hyperref[def:formal-difference-equivalence]{Formal-Difference Equivalence}
  \item \hyperref[lem:positivity-respects-equivalence]{Positivity Respects Formal-Difference Equivalence}
\end{itemize}
\end{dependencies}

\begin{definitionbox}{Definition}
\begin{definition}[Strict Order on $\mathbb{Z}$]
\label{def:strict-order-on-z}
For $x,y\in\mathbb{Z}$, define
\[
x<y \quad:\Longleftrightarrow\quad y-x \text{ is positive}.
\]
\end{definition}
\end{definitionbox}

\begin{remark*}[Standard quantified statement]
\[
\forall x,y\in\mathbb{Z},\quad
x<_{\mathbb{Z}}y
\Longleftrightarrow
\operatorname{Positive}_{\mathbb{Z}}(y-x).
\]
\end{remark*}

\begin{remark*}[Predicate reading]
\[
\operatorname{StrictOrderFromPositiveCone}(<_{\mathbb{Z}},\operatorname{Positive}_{\mathbb{Z}},\mathbb{Z}).
\]
\end{remark*}

\begin{remark*}[Interpretation]
Strict order is induced by positivity of the difference.
\end{remark*}

\begin{dependencies}
\begin{itemize}
  \item \hyperref[def:subtraction-on-z]{Subtraction on $\mathbb{Z}$}
  \item \hyperref[def:positive-integer]{Positive Integer}
\end{itemize}
\end{dependencies}

\begin{definitionbox}{Definition}
\begin{definition}[Order on $\mathbb{Z}$]
\label{def:order-on-z}
For $x,y\in\mathbb{Z}$, define
\[
x\le y \quad:\Longleftrightarrow\quad x<y \text{ or } x=y.
\]
\end{definition}
\end{definitionbox}

\begin{remark*}[Standard quantified statement]
\[
\forall x,y\in\mathbb{Z},\quad
x\le_{\mathbb{Z}}y
\Longleftrightarrow
x<_{\mathbb{Z}}y\lor x=y.
\]
\end{remark*}

\begin{remark*}[Predicate reading]
\[
\operatorname{NonStrictOrderFromStrictOrder}(\le_{\mathbb{Z}},<_{\mathbb{Z}},\mathbb{Z}).
\]
\end{remark*}

\begin{remark*}[Interpretation]
The non-strict order is obtained by adjoining equality to the strict order.
\end{remark*}

\begin{dependencies}
\begin{itemize}
  \item \hyperref[def:strict-order-on-z]{Strict Order on $\mathbb{Z}$}
\end{itemize}
\end{dependencies}

\begin{propositionbox}{Proposition}
\begin{proposition}[Product of Positives Is Positive]
\label{prop:product-of-positives-is-positive}
For all $x,y\in\mathbb{Z}$,
\[
0<x\land 0<y \Longrightarrow 0<x\cdot y.
\]
\hyperref[prf:product-of-positives-is-positive]{\textit{Go to proof.}}
\end{proposition}
\end{propositionbox}

\begin{remark*}[Standard quantified statement]
\[
\forall x,y\in\mathbb{Z},\quad
0<x\land 0<y \Longrightarrow 0<x\cdot y.
\]
\end{remark*}

\begin{remark*}[Predicate reading]
\[
\operatorname{PositiveConeMultiplicativelyClosed}(\operatorname{Positive}_{\mathbb{Z}},\mathbb{Z},\cdot_{\mathbb{Z}}).
\]
\end{remark*}

\begin{remark*}[Interpretation]
The positive integers are closed under multiplication.
\end{remark*}

\begin{dependencies}
\begin{itemize}
  \item \hyperref[def:positive-integer]{Positive Integer}
  \item \hyperref[def:multiplication-on-z-classes]{Multiplication on $\mathbb{Z}$}
  \item \hyperref[thm:multiplication-on-z-well-defined]{Multiplication on $\mathbb{Z}$ Is Well-Defined}
  \item \hyperref[def:order-on-w]{Order on $\mathbb{W}$}
\end{itemize}
\end{dependencies}

\subsection*{Canonical Form and Sign}

\begin{theorembox}{Theorem}
\begin{theorem}[Canonical Form of Integers]
\label{thm:canonical-form-of-integers}
For every $x\in\mathbb{Z}$, exactly one of the following holds:
\[
x=[n,0]\text{ for some nonzero }n\in\mathbb{W},\qquad
x=[0,0],\qquad
x=[0,n]\text{ for some nonzero }n\in\mathbb{W}.
\]
\hyperref[prf:canonical-form-of-integers]{\textit{Go to proof.}}
\end{theorem}
\end{theorembox}

\begin{remark*}[Standard quantified statement]
\[
\forall x\in\mathbb{Z},\quad
\text{exactly one of }x=[n,0]\ (n\ne0),\ x=[0,0],\ x=[0,n]\ (n\ne0)\text{ holds}.
\]
\end{remark*}

\begin{remark*}[Predicate reading]
\[
\operatorname{IntegerCanonicalTrichotomy}(\mathbb{Z},0_{\mathbb{Z}}).
\]
\end{remark*}

\begin{remark*}[Interpretation]
Every integer has a unique positive, zero, or negative representative form.
\end{remark*}

\begin{dependencies}
\begin{itemize}
  \item \hyperref[def:integers]{Integers}
  \item \hyperref[def:integer-class-notation]{Integer Class Notation}
  \item \hyperref[def:formal-difference-equivalence]{Formal-Difference Equivalence}
  \item \hyperref[def:whole-numbers]{Whole Numbers}
  \item \hyperref[thm:order-on-w-total]{Order on $\mathbb{W}$ Is Total}
\end{itemize}
\end{dependencies}

\begin{corollarybox}{Corollary}
\begin{corollary}[Integers Split by Sign]
\label{cor:integers-split-by-sign}
Every integer is nonnegative or the negative of a nonnegative integer, and
the two descriptions overlap only at $0_{\mathbb{Z}}$.
\hyperref[prf:integers-split-by-sign]{\textit{Go to proof.}}
\end{corollary}
\end{corollarybox}

\begin{remark*}[Standard quantified statement]
\[
\forall x\in\mathbb{Z},\quad
0_{\mathbb{Z}}\le x
\quad\text{or}\quad
\exists n\in\mathbb{W},\ x=-\iota(n),
\]
with overlap only at $0_{\mathbb{Z}}$.
\end{remark*}

\begin{remark*}[Predicate reading]
\[
\operatorname{SignDecomposition}(\mathbb{Z},\le_{\mathbb{Z}},0_{\mathbb{Z}}).
\]
\end{remark*}

\begin{remark*}[Interpretation]
The canonical form splits integers into nonnegative and negative parts.
\end{remark*}

\begin{dependencies}
\begin{itemize}
  \item \hyperref[thm:canonical-form-of-integers]{Canonical Form of Integers}
  \item \hyperref[def:positive-integer]{Positive Integer}
  \item \hyperref[def:zero-in-z]{Zero in $\mathbb{Z}$}
\end{itemize}
\end{dependencies}

\begin{corollarybox}{Corollary}
\begin{corollary}[Image of the Embedding Is the Nonnegative Integers]
\label{cor:image-of-embedding-is-nonnegatives}
The image of the canonical embedding $\mathbb{W}\to\mathbb{Z}$ is exactly
the set of nonnegative integers.
\hyperref[prf:image-of-embedding-is-nonnegatives]{\textit{Go to proof.}}
\end{corollary}
\end{corollarybox}

\begin{remark*}[Standard quantified statement]
\[
\operatorname{im}(\iota)
=
\{x\in\mathbb{Z}:0_{\mathbb{Z}}\le x\}.
\]
\end{remark*}

\begin{remark*}[Predicate reading]
\[
\operatorname{ImageEqualsNonnegativeCone}(\iota,\mathbb{W},\mathbb{Z}).
\]
\end{remark*}

\begin{remark*}[Interpretation]
The embedded whole numbers are exactly the nonnegative integers.
\end{remark*}

\begin{dependencies}
\begin{itemize}
  \item \hyperref[def:canonical-embedding-w-into-z]{Embedding of $\mathbb{W}$ into $\mathbb{Z}$}
  \item \hyperref[cor:integers-split-by-sign]{Integers Split by Sign}
  \item \hyperref[thm:embedding-w-into-z-injective]{Embedding $\mathbb{W}$ into $\mathbb{Z}$ Is Injective}
\end{itemize}
\end{dependencies}

\begin{theorembox}{Theorem}
\begin{theorem}[Sign Trichotomy on $\mathbb{Z}$]
\label{thm:sign-trichotomy-on-z}
For every $x\in\mathbb{Z}$, exactly one of
\[
x<0,\qquad x=0,\qquad 0<x
\]
holds.
\hyperref[prf:sign-trichotomy-on-z]{\textit{Go to proof.}}
\end{theorem}
\end{theorembox}

\begin{remark*}[Standard quantified statement]
\[
\forall x\in\mathbb{Z},\quad
\text{exactly one of }x<0,\ x=0,\ 0<x\text{ holds}.
\]
\end{remark*}

\begin{remark*}[Predicate reading]
\[
\operatorname{SignTrichotomy}(\mathbb{Z},<_{\mathbb{Z}},0_{\mathbb{Z}}).
\]
\end{remark*}

\begin{remark*}[Interpretation]
Every integer has exactly one sign relative to zero.
\end{remark*}

\begin{dependencies}
\begin{itemize}
  \item \hyperref[thm:canonical-form-of-integers]{Canonical Form of Integers}
  \item \hyperref[thm:order-on-w-total]{Order on $\mathbb{W}$ Is Total}
\end{itemize}
\end{dependencies}

\begin{theorembox}{Theorem}
\begin{theorem}[Trichotomy on $\mathbb{Z}$]
\label{thm:trichotomy-on-z}
For all $x,y\in\mathbb{Z}$, exactly one of
\[
x<y,\qquad x=y,\qquad y<x
\]
holds.
\hyperref[prf:trichotomy-on-z]{\textit{Go to proof.}}
\end{theorem}
\end{theorembox}

\begin{remark*}[Standard quantified statement]
\[
\forall x,y\in\mathbb{Z},\quad
\text{exactly one of }x<y,\ x=y,\ y<x\text{ holds}.
\]
\end{remark*}

\begin{remark*}[Predicate reading]
\[
\operatorname{Trichotomous}(<_{\mathbb{Z}},\mathbb{Z}).
\]
\end{remark*}

\begin{remark*}[Interpretation]
The strict order compares every pair of integers in exactly one way.
\end{remark*}

\begin{dependencies}
\begin{itemize}
  \item \hyperref[def:strict-order-on-z]{Strict Order on $\mathbb{Z}$}
  \item \hyperref[def:subtraction-on-z]{Subtraction on $\mathbb{Z}$}
  \item \hyperref[thm:sign-trichotomy-on-z]{Sign Trichotomy on $\mathbb{Z}$}
\end{itemize}
\end{dependencies}

\subsection*{\texorpdfstring{Total Order on $\mathbb{Z}$}{Total Order on Z}}

\begin{theorembox}{Theorem}
\begin{theorem}[Strict Order on $\mathbb{Z}$ Is Irreflexive]
\label{thm:strict-order-on-z-irreflexive}
For every $x\in\mathbb{Z}$, not $x<x$.
\hyperref[prf:strict-order-on-z-irreflexive]{\textit{Go to proof.}}
\end{theorem}
\end{theorembox}

\begin{remark*}[Standard quantified statement]
\[
\forall x\in\mathbb{Z},\quad \neg(x<_{\mathbb{Z}}x).
\]
\end{remark*}

\begin{remark*}[Predicate reading]
\[
\operatorname{Irreflexive}(<_{\mathbb{Z}},\mathbb{Z}).
\]
\end{remark*}

\begin{remark*}[Interpretation]
No integer is strictly less than itself.
\end{remark*}

\begin{dependencies}
\begin{itemize}
  \item \hyperref[def:strict-order-on-z]{Strict Order on $\mathbb{Z}$}
  \item \hyperref[def:subtraction-on-z]{Subtraction on $\mathbb{Z}$}
  \item \hyperref[thm:zero-is-an-integer]{Zero Is an Integer}
\end{itemize}
\end{dependencies}

\begin{theorembox}{Theorem}
\begin{theorem}[Strict Order on $\mathbb{Z}$ Is Asymmetric]
\label{thm:strict-order-on-z-asymmetric}
For all $x,y\in\mathbb{Z}$,
\[
x<y \Longrightarrow \neg(y<x).
\]
\hyperref[prf:strict-order-on-z-asymmetric]{\textit{Go to proof.}}
\end{theorem}
\end{theorembox}

\begin{remark*}[Standard quantified statement]
\[
\forall x,y\in\mathbb{Z},\quad
x<_{\mathbb{Z}}y \Longrightarrow \neg(y<_{\mathbb{Z}}x).
\]
\end{remark*}

\begin{remark*}[Predicate reading]
\[
\operatorname{Asymmetric}(<_{\mathbb{Z}},\mathbb{Z}).
\]
\end{remark*}

\begin{remark*}[Interpretation]
Strict order cannot hold in both directions.
\end{remark*}

\begin{dependencies}
\begin{itemize}
  \item \hyperref[def:strict-order-on-z]{Strict Order on $\mathbb{Z}$}
  \item \hyperref[thm:strict-order-on-z-irreflexive]{Strict Order on $\mathbb{Z}$ Is Irreflexive}
  \item \hyperref[thm:trichotomy-on-z]{Trichotomy on $\mathbb{Z}$}
\end{itemize}
\end{dependencies}

\begin{theorembox}{Theorem}
\begin{theorem}[Strict Order on $\mathbb{Z}$ Is Transitive]
\label{thm:strict-order-on-z-transitive}
For all $x,y,z\in\mathbb{Z}$,
\[
x<y\land y<z \Longrightarrow x<z.
\]
\hyperref[prf:strict-order-on-z-transitive]{\textit{Go to proof.}}
\end{theorem}
\end{theorembox}

\begin{remark*}[Standard quantified statement]
\[
\forall x,y,z\in\mathbb{Z},\quad
x<y\land y<z \Longrightarrow x<z.
\]
\end{remark*}

\begin{remark*}[Predicate reading]
\[
\operatorname{Transitive}(<_{\mathbb{Z}},\mathbb{Z}).
\]
\end{remark*}

\begin{remark*}[Interpretation]
Strict integer inequalities compose transitively.
\end{remark*}

\begin{dependencies}
\begin{itemize}
  \item \hyperref[def:strict-order-on-z]{Strict Order on $\mathbb{Z}$}
  \item \hyperref[def:subtraction-on-z]{Subtraction on $\mathbb{Z}$}
  \item \hyperref[def:positive-integer]{Positive Integer}
  \item \hyperref[thm:addition-preserves-strict-order-on-z]{Addition Preserves Strict Order on $\mathbb{Z}$}
\end{itemize}
\end{dependencies}

\begin{theorembox}{Theorem}
\begin{theorem}[Order on $\mathbb{Z}$ Is Reflexive]
\label{thm:order-on-z-reflexive}
For every $x\in\mathbb{Z}$, $x\le x$.
\hyperref[prf:order-on-z-reflexive]{\textit{Go to proof.}}
\end{theorem}
\end{theorembox}

\begin{remark*}[Standard quantified statement]
\[
\forall x\in\mathbb{Z},\quad x\le_{\mathbb{Z}}x.
\]
\end{remark*}

\begin{remark*}[Predicate reading]
\[
\operatorname{Reflexive}(\le_{\mathbb{Z}},\mathbb{Z}).
\]
\end{remark*}

\begin{remark*}[Interpretation]
Every integer is less than or equal to itself.
\end{remark*}

\begin{dependencies}
\begin{itemize}
  \item \hyperref[def:order-on-z]{Order on $\mathbb{Z}$}
  \item \hyperref[def:strict-order-on-z]{Strict Order on $\mathbb{Z}$}
\end{itemize}
\end{dependencies}

\begin{theorembox}{Theorem}
\begin{theorem}[Order on $\mathbb{Z}$ Is Antisymmetric]
\label{thm:order-on-z-antisymmetric}
For all $x,y\in\mathbb{Z}$,
\[
x\le y\land y\le x \Longrightarrow x=y.
\]
\hyperref[prf:order-on-z-antisymmetric]{\textit{Go to proof.}}
\end{theorem}
\end{theorembox}

\begin{remark*}[Standard quantified statement]
\[
\forall x,y\in\mathbb{Z},\quad
x\le y\land y\le x \Longrightarrow x=y.
\]
\end{remark*}

\begin{remark*}[Predicate reading]
\[
\operatorname{Antisymmetric}(\le_{\mathbb{Z}},\mathbb{Z}).
\]
\end{remark*}

\begin{remark*}[Interpretation]
Mutual non-strict comparison forces equality.
\end{remark*}

\begin{dependencies}
\begin{itemize}
  \item \hyperref[def:order-on-z]{Order on $\mathbb{Z}$}
  \item \hyperref[def:strict-order-on-z]{Strict Order on $\mathbb{Z}$}
  \item \hyperref[thm:trichotomy-on-z]{Trichotomy on $\mathbb{Z}$}
\end{itemize}
\end{dependencies}

\begin{theorembox}{Theorem}
\begin{theorem}[Order on $\mathbb{Z}$ Is Transitive]
\label{thm:order-on-z-transitive}
For all $x,y,z\in\mathbb{Z}$,
\[
x\le y\land y\le z \Longrightarrow x\le z.
\]
\hyperref[prf:order-on-z-transitive]{\textit{Go to proof.}}
\end{theorem}
\end{theorembox}

\begin{remark*}[Standard quantified statement]
\[
\forall x,y,z\in\mathbb{Z},\quad
x\le y\land y\le z \Longrightarrow x\le z.
\]
\end{remark*}

\begin{remark*}[Predicate reading]
\[
\operatorname{Transitive}(\le_{\mathbb{Z}},\mathbb{Z}).
\]
\end{remark*}

\begin{remark*}[Interpretation]
Non-strict integer inequalities compose transitively.
\end{remark*}

\begin{dependencies}
\begin{itemize}
  \item \hyperref[def:order-on-z]{Order on $\mathbb{Z}$}
  \item \hyperref[def:strict-order-on-z]{Strict Order on $\mathbb{Z}$}
  \item \hyperref[thm:strict-order-on-z-transitive]{Strict Order on $\mathbb{Z}$ Is Transitive}
  \item \hyperref[thm:trichotomy-on-z]{Trichotomy on $\mathbb{Z}$}
\end{itemize}
\end{dependencies}

\begin{theorembox}{Theorem}
\begin{theorem}[Order on $\mathbb{Z}$ Is Total]
\label{thm:order-on-z-total}
For all $x,y\in\mathbb{Z}$,
\[
x\le y\text{ or }y\le x.
\]
\hyperref[prf:order-on-z-total]{\textit{Go to proof.}}
\end{theorem}
\end{theorembox}

\begin{remark*}[Standard quantified statement]
\[
\forall x,y\in\mathbb{Z},\quad x\le y\lor y\le x.
\]
\end{remark*}

\begin{remark*}[Predicate reading]
\[
\operatorname{Total}(\le_{\mathbb{Z}},\mathbb{Z}).
\]
\end{remark*}

\begin{remark*}[Interpretation]
Every two integers are comparable in the non-strict order.
\end{remark*}

\begin{dependencies}
\begin{itemize}
  \item \hyperref[def:order-on-z]{Order on $\mathbb{Z}$}
  \item \hyperref[thm:trichotomy-on-z]{Trichotomy on $\mathbb{Z}$}
\end{itemize}
\end{dependencies}

\begin{theorembox}{Theorem}
\begin{theorem}[$\mathbb{Z}$ Is a Totally Ordered Set]
\label{thm:z-totally-ordered-set}
The ordered pair $(\mathbb{Z},\le)$ is a totally ordered set.
\hyperref[prf:z-totally-ordered-set]{\textit{Go to proof.}}
\end{theorem}
\end{theorembox}

\begin{remark*}[Standard quantified statement]
\[
\operatorname{Reflexive}(\le_{\mathbb{Z}},\mathbb{Z})
\land
\operatorname{Antisymmetric}(\le_{\mathbb{Z}},\mathbb{Z})
\land
\operatorname{Transitive}(\le_{\mathbb{Z}},\mathbb{Z})
\land
\operatorname{Total}(\le_{\mathbb{Z}},\mathbb{Z}).
\]
\end{remark*}

\begin{remark*}[Predicate reading]
\[
\operatorname{TotalOrder}(\le_{\mathbb{Z}},\mathbb{Z}).
\]
\end{remark*}

\begin{remark*}[Interpretation]
The non-strict integer order is a total order.
\end{remark*}

\begin{dependencies}
\begin{itemize}
  \item \hyperref[def:ordered-set]{Ordered Set}
  \item \hyperref[def:total-order]{Total Order}
  \item \hyperref[def:order-on-z]{Order on $\mathbb{Z}$}
  \item \hyperref[thm:order-on-z-reflexive]{Order on $\mathbb{Z}$ Is Reflexive}
  \item \hyperref[thm:order-on-z-antisymmetric]{Order on $\mathbb{Z}$ Is Antisymmetric}
  \item \hyperref[thm:order-on-z-transitive]{Order on $\mathbb{Z}$ Is Transitive}
  \item \hyperref[thm:order-on-z-total]{Order on $\mathbb{Z}$ Is Total}
\end{itemize}
\end{dependencies}

\subsection*{\texorpdfstring{Addition Order Compatibility on $\mathbb{Z}$}{Addition Order Compatibility on Z}}

\begin{theorembox}{Theorem}
\begin{theorem}[Left Addition Preserves Strict Order on $\mathbb{Z}$]
\label{thm:left-addition-preserves-strict-order-on-z}
For all $x,y,z\in\mathbb{Z}$,
\[
x<y \Longrightarrow z+x<z+y.
\]
\hyperref[prf:left-addition-preserves-strict-order-on-z]{\textit{Go to proof.}}
\end{theorem}
\end{theorembox}

\begin{remark*}[Standard quantified statement]
\[
\forall x,y,z\in\mathbb{Z},\quad
x<y \Longrightarrow z+x<z+y.
\]
\end{remark*}

\begin{remark*}[Predicate reading]
\[
\operatorname{RightMonotone}(+_{\mathbb{Z}},<_{\mathbb{Z}},\mathbb{Z}).
\]
\end{remark*}

\begin{remark*}[Interpretation]
Adding a fixed integer on the left preserves strict order.
\end{remark*}

\begin{dependencies}
\begin{itemize}
  \item \hyperref[def:strict-order-on-z]{Strict Order on $\mathbb{Z}$}
  \item \hyperref[def:addition-on-z-classes]{Addition on $\mathbb{Z}$}
  \item \hyperref[thm:addition-on-z-well-defined]{Addition on $\mathbb{Z}$ Is Well-Defined}
\end{itemize}
\end{dependencies}

\begin{corollarybox}{Corollary}
\begin{corollary}[Right Addition Preserves Strict Order on $\mathbb{Z}$]
\label{cor:right-addition-preserves-strict-order-on-z}
For all $x,y,z\in\mathbb{Z}$,
\[
x<y \Longrightarrow x+z<y+z.
\]
\hyperref[prf:right-addition-preserves-strict-order-on-z]{\textit{Go to proof.}}
\end{corollary}
\end{corollarybox}

\begin{remark*}[Standard quantified statement]
\[
\forall x,y,z\in\mathbb{Z},\quad
x<y \Longrightarrow x+z<y+z.
\]
\end{remark*}

\begin{remark*}[Predicate reading]
\[
\operatorname{LeftMonotone}(+_{\mathbb{Z}},<_{\mathbb{Z}},\mathbb{Z}).
\]
\end{remark*}

\begin{remark*}[Interpretation]
Adding a fixed integer on the right preserves strict order.
\end{remark*}

\begin{dependencies}
\begin{itemize}
  \item \hyperref[thm:left-addition-preserves-strict-order-on-z]{Left Addition Preserves Strict Order on $\mathbb{Z}$}
  \item \hyperref[thm:addition-on-z-commutative]{Addition on $\mathbb{Z}$ Is Commutative}
\end{itemize}
\end{dependencies}

\begin{theorembox}{Theorem}
\begin{theorem}[Addition Preserves Strict Order on $\mathbb{Z}$]
\label{thm:addition-preserves-strict-order-on-z}
Addition preserves strict order on both sides in $\mathbb{Z}$.
\hyperref[prf:addition-preserves-strict-order-on-z]{\textit{Go to proof.}}
\end{theorem}
\end{theorembox}

\begin{remark*}[Standard quantified statement]
\[
\operatorname{LeftMonotone}(+_{\mathbb{Z}},<_{\mathbb{Z}},\mathbb{Z})
\quad\text{and}\quad
\operatorname{RightMonotone}(+_{\mathbb{Z}},<_{\mathbb{Z}},\mathbb{Z}).
\]
\end{remark*}

\begin{remark*}[Predicate reading]
\[
\operatorname{MonotoneOperation}(+_{\mathbb{Z}},<_{\mathbb{Z}},\mathbb{Z}).
\]
\end{remark*}

\begin{remark*}[Interpretation]
Addition preserves strict order in both operation slots.
\end{remark*}

\begin{dependencies}
\begin{itemize}
  \item \hyperref[thm:left-addition-preserves-strict-order-on-z]{Left Addition Preserves Strict Order on $\mathbb{Z}$}
  \item \hyperref[cor:right-addition-preserves-strict-order-on-z]{Right Addition Preserves Strict Order on $\mathbb{Z}$}
\end{itemize}
\end{dependencies}

\begin{theorembox}{Theorem}
\begin{theorem}[Left Addition Preserves Order on $\mathbb{Z}$]
\label{thm:left-addition-preserves-order-on-z}
For all $x,y,z\in\mathbb{Z}$,
\[
x\le y \Longrightarrow z+x\le z+y.
\]
\hyperref[prf:left-addition-preserves-order-on-z]{\textit{Go to proof.}}
\end{theorem}
\end{theorembox}

\begin{remark*}[Standard quantified statement]
\[
\forall x,y,z\in\mathbb{Z},\quad
x\le y \Longrightarrow z+x\le z+y.
\]
\end{remark*}

\begin{remark*}[Predicate reading]
\[
\operatorname{RightMonotone}(+_{\mathbb{Z}},\le_{\mathbb{Z}},\mathbb{Z}).
\]
\end{remark*}

\begin{remark*}[Interpretation]
Adding a fixed integer on the left preserves non-strict order.
\end{remark*}

\begin{dependencies}
\begin{itemize}
  \item \hyperref[def:order-on-z]{Order on $\mathbb{Z}$}
  \item \hyperref[thm:left-addition-preserves-strict-order-on-z]{Left Addition Preserves Strict Order on $\mathbb{Z}$}
  \item \hyperref[thm:addition-on-z-well-defined]{Addition on $\mathbb{Z}$ Is Well-Defined}
\end{itemize}
\end{dependencies}

\begin{corollarybox}{Corollary}
\begin{corollary}[Right Addition Preserves Order on $\mathbb{Z}$]
\label{cor:right-addition-preserves-order-on-z}
For all $x,y,z\in\mathbb{Z}$,
\[
x\le y \Longrightarrow x+z\le y+z.
\]
\hyperref[prf:right-addition-preserves-order-on-z]{\textit{Go to proof.}}
\end{corollary}
\end{corollarybox}

\begin{remark*}[Standard quantified statement]
\[
\forall x,y,z\in\mathbb{Z},\quad
x\le y \Longrightarrow x+z\le y+z.
\]
\end{remark*}

\begin{remark*}[Predicate reading]
\[
\operatorname{LeftMonotone}(+_{\mathbb{Z}},\le_{\mathbb{Z}},\mathbb{Z}).
\]
\end{remark*}

\begin{remark*}[Interpretation]
Adding a fixed integer on the right preserves non-strict order.
\end{remark*}

\begin{dependencies}
\begin{itemize}
  \item \hyperref[thm:left-addition-preserves-order-on-z]{Left Addition Preserves Order on $\mathbb{Z}$}
  \item \hyperref[thm:addition-on-z-commutative]{Addition on $\mathbb{Z}$ Is Commutative}
\end{itemize}
\end{dependencies}

\begin{theorembox}{Theorem}
\begin{theorem}[Addition Preserves Order on $\mathbb{Z}$]
\label{thm:addition-preserves-order-on-z}
Addition preserves non-strict order on both sides in $\mathbb{Z}$.
\hyperref[prf:addition-preserves-order-on-z]{\textit{Go to proof.}}
\end{theorem}
\end{theorembox}

\begin{remark*}[Standard quantified statement]
\[
\operatorname{LeftMonotone}(+_{\mathbb{Z}},\le_{\mathbb{Z}},\mathbb{Z})
\quad\text{and}\quad
\operatorname{RightMonotone}(+_{\mathbb{Z}},\le_{\mathbb{Z}},\mathbb{Z}).
\]
\end{remark*}

\begin{remark*}[Predicate reading]
\[
\operatorname{MonotoneOperation}(+_{\mathbb{Z}},\le_{\mathbb{Z}},\mathbb{Z}).
\]
\end{remark*}

\begin{remark*}[Interpretation]
Addition preserves non-strict order in both operation slots.
\end{remark*}

\begin{dependencies}
\begin{itemize}
  \item \hyperref[thm:left-addition-preserves-order-on-z]{Left Addition Preserves Order on $\mathbb{Z}$}
  \item \hyperref[cor:right-addition-preserves-order-on-z]{Right Addition Preserves Order on $\mathbb{Z}$}
\end{itemize}
\end{dependencies}

\begin{corollarybox}{Corollary}
\begin{corollary}[Addition Reflects Order on $\mathbb{Z}$]
\label{cor:addition-reflects-order-on-z}
Addition by a fixed integer reflects both strict and non-strict order.
\hyperref[prf:addition-reflects-order-on-z]{\textit{Go to proof.}}
\end{corollary}
\end{corollarybox}

\begin{remark*}[Standard quantified statement]
\[
\text{For each }z\in\mathbb{Z},\quad
x+z\le y+z \Longrightarrow x\le y,
\]
and the analogous strict-order implication holds.
\end{remark*}

\begin{remark*}[Predicate reading]
\[
\operatorname{OrderReflectingOperation}(+_{\mathbb{Z}},\le_{\mathbb{Z}},\mathbb{Z})
\quad\text{and}\quad
\operatorname{OrderReflectingOperation}(+_{\mathbb{Z}},<_{\mathbb{Z}},\mathbb{Z}).
\]
\end{remark*}

\begin{remark*}[Interpretation]
Additive inverses allow cancellation of a fixed summand in order comparisons.
\end{remark*}

\begin{dependencies}
\begin{itemize}
  \item \hyperref[thm:addition-preserves-strict-order-on-z]{Addition Preserves Strict Order on $\mathbb{Z}$}
  \item \hyperref[thm:addition-preserves-order-on-z]{Addition Preserves Order on $\mathbb{Z}$}
  \item \hyperref[thm:negation-left-additive-inverse-on-z]{Negation Is a Left Additive Inverse on $\mathbb{Z}$}
\end{itemize}
\end{dependencies}

\subsection*{\texorpdfstring{Multiplication Order Compatibility on $\mathbb{Z}$}{Multiplication Order Compatibility on Z}}

\begin{theorembox}{Theorem}
\begin{theorem}[Left Multiplication by Positives Preserves Strict Order on $\mathbb{Z}$]
\label{thm:left-positive-multiplication-preserves-strict-order-on-z}
For all $x,y,z\in\mathbb{Z}$,
\[
0<z\land x<y \Longrightarrow z\cdot x<z\cdot y.
\]
\hyperref[prf:left-positive-multiplication-preserves-strict-order-on-z]{\textit{Go to proof.}}
\end{theorem}
\end{theorembox}

\begin{remark*}[Standard quantified statement]
\[
\forall x,y,z\in\mathbb{Z},\quad
0<z\land x<y \Longrightarrow z\cdot x<z\cdot y.
\]
\end{remark*}

\begin{remark*}[Predicate reading]
\[
\operatorname{PositiveLeftFactorPreservesOrder}(\cdot_{\mathbb{Z}},<_{\mathbb{Z}},0_{\mathbb{Z}},\mathbb{Z}).
\]
\end{remark*}

\begin{remark*}[Interpretation]
Multiplication by a positive left factor preserves strict order.
\end{remark*}

\begin{dependencies}
\begin{itemize}
  \item \hyperref[def:strict-order-on-z]{Strict Order on $\mathbb{Z}$}
  \item \hyperref[prop:product-of-positives-is-positive]{Product of Positives Is Positive}
  \item \hyperref[thm:multiplication-on-z-well-defined]{Multiplication on $\mathbb{Z}$ Is Well-Defined}
  \item \hyperref[thm:multiplication-distributes-over-addition-on-z]{Multiplication Distributes over Addition on $\mathbb{Z}$}
\end{itemize}
\end{dependencies}

\begin{corollarybox}{Corollary}
\begin{corollary}[Right Multiplication by Positives Preserves Strict Order on $\mathbb{Z}$]
\label{cor:right-positive-multiplication-preserves-strict-order-on-z}
For all $x,y,z\in\mathbb{Z}$,
\[
0<z\land x<y \Longrightarrow x\cdot z<y\cdot z.
\]
\hyperref[prf:right-positive-multiplication-preserves-strict-order-on-z]{\textit{Go to proof.}}
\end{corollary}
\end{corollarybox}

\begin{remark*}[Standard quantified statement]
\[
\forall x,y,z\in\mathbb{Z},\quad
0<z\land x<y \Longrightarrow x\cdot z<y\cdot z.
\]
\end{remark*}

\begin{remark*}[Predicate reading]
\[
\operatorname{PositiveRightFactorPreservesOrder}(\cdot_{\mathbb{Z}},<_{\mathbb{Z}},0_{\mathbb{Z}},\mathbb{Z}).
\]
\end{remark*}

\begin{remark*}[Interpretation]
Multiplication by a positive right factor preserves strict order.
\end{remark*}

\begin{dependencies}
\begin{itemize}
  \item \hyperref[thm:left-positive-multiplication-preserves-strict-order-on-z]{Left Multiplication by Positives Preserves Strict Order on $\mathbb{Z}$}
  \item \hyperref[thm:multiplication-on-z-commutative]{Multiplication on $\mathbb{Z}$ Is Commutative}
\end{itemize}
\end{dependencies}

\begin{theorembox}{Theorem}
\begin{theorem}[Multiplication by Positives Preserves Strict Order on $\mathbb{Z}$]
\label{thm:positive-multiplication-preserves-strict-order-on-z}
Multiplication by a positive integer preserves strict order on both sides.
\hyperref[prf:positive-multiplication-preserves-strict-order-on-z]{\textit{Go to proof.}}
\end{theorem}
\end{theorembox}

\begin{remark*}[Standard quantified statement]
\[
\operatorname{PositiveLeftFactorPreservesOrder}(\cdot_{\mathbb{Z}},<_{\mathbb{Z}},0_{\mathbb{Z}},\mathbb{Z})
\quad\text{and}\quad
\operatorname{PositiveRightFactorPreservesOrder}(\cdot_{\mathbb{Z}},<_{\mathbb{Z}},0_{\mathbb{Z}},\mathbb{Z}).
\]
\end{remark*}

\begin{remark*}[Predicate reading]
\[
\operatorname{PositiveFactorPreservesOrder}(\cdot_{\mathbb{Z}},<_{\mathbb{Z}},0_{\mathbb{Z}},\mathbb{Z}).
\]
\end{remark*}

\begin{remark*}[Interpretation]
Multiplication by positive integers preserves strict order on both sides.
\end{remark*}

\begin{dependencies}
\begin{itemize}
  \item \hyperref[thm:left-positive-multiplication-preserves-strict-order-on-z]{Left Multiplication by Positives Preserves Strict Order on $\mathbb{Z}$}
  \item \hyperref[cor:right-positive-multiplication-preserves-strict-order-on-z]{Right Multiplication by Positives Preserves Strict Order on $\mathbb{Z}$}
\end{itemize}
\end{dependencies}

\begin{theorembox}{Theorem}
\begin{theorem}[Left Multiplication by Positives Preserves Order on $\mathbb{Z}$]
\label{thm:left-positive-multiplication-preserves-order-on-z}
For all $x,y,z\in\mathbb{Z}$,
\[
0<z\land x\le y \Longrightarrow z\cdot x\le z\cdot y.
\]
\hyperref[prf:left-positive-multiplication-preserves-order-on-z]{\textit{Go to proof.}}
\end{theorem}
\end{theorembox}

\begin{remark*}[Standard quantified statement]
\[
\forall x,y,z\in\mathbb{Z},\quad
0<z\land x\le y \Longrightarrow z\cdot x\le z\cdot y.
\]
\end{remark*}

\begin{remark*}[Predicate reading]
\[
\operatorname{PositiveLeftFactorPreservesOrder}(\cdot_{\mathbb{Z}},\le_{\mathbb{Z}},0_{\mathbb{Z}},\mathbb{Z}).
\]
\end{remark*}

\begin{remark*}[Interpretation]
Multiplication by a positive left factor preserves non-strict order.
\end{remark*}

\begin{dependencies}
\begin{itemize}
  \item \hyperref[def:order-on-z]{Order on $\mathbb{Z}$}
  \item \hyperref[thm:left-positive-multiplication-preserves-strict-order-on-z]{Left Multiplication by Positives Preserves Strict Order on $\mathbb{Z}$}
\end{itemize}
\end{dependencies}

\begin{corollarybox}{Corollary}
\begin{corollary}[Right Multiplication by Positives Preserves Order on $\mathbb{Z}$]
\label{cor:right-positive-multiplication-preserves-order-on-z}
For all $x,y,z\in\mathbb{Z}$,
\[
0<z\land x\le y \Longrightarrow x\cdot z\le y\cdot z.
\]
\hyperref[prf:right-positive-multiplication-preserves-order-on-z]{\textit{Go to proof.}}
\end{corollary}
\end{corollarybox}

\begin{remark*}[Standard quantified statement]
\[
\forall x,y,z\in\mathbb{Z},\quad
0<z\land x\le y \Longrightarrow x\cdot z\le y\cdot z.
\]
\end{remark*}

\begin{remark*}[Predicate reading]
\[
\operatorname{PositiveRightFactorPreservesOrder}(\cdot_{\mathbb{Z}},\le_{\mathbb{Z}},0_{\mathbb{Z}},\mathbb{Z}).
\]
\end{remark*}

\begin{remark*}[Interpretation]
Multiplication by a positive right factor preserves non-strict order.
\end{remark*}

\begin{dependencies}
\begin{itemize}
  \item \hyperref[thm:left-positive-multiplication-preserves-order-on-z]{Left Multiplication by Positives Preserves Order on $\mathbb{Z}$}
  \item \hyperref[thm:multiplication-on-z-commutative]{Multiplication on $\mathbb{Z}$ Is Commutative}
\end{itemize}
\end{dependencies}

\begin{theorembox}{Theorem}
\begin{theorem}[Multiplication by Positives Preserves Order on $\mathbb{Z}$]
\label{thm:positive-multiplication-preserves-order-on-z}
Multiplication by a positive integer preserves non-strict order on both sides.
\hyperref[prf:positive-multiplication-preserves-order-on-z]{\textit{Go to proof.}}
\end{theorem}
\end{theorembox}

\begin{remark*}[Standard quantified statement]
\[
\operatorname{PositiveLeftFactorPreservesOrder}(\cdot_{\mathbb{Z}},\le_{\mathbb{Z}},0_{\mathbb{Z}},\mathbb{Z})
\quad\text{and}\quad
\operatorname{PositiveRightFactorPreservesOrder}(\cdot_{\mathbb{Z}},\le_{\mathbb{Z}},0_{\mathbb{Z}},\mathbb{Z}).
\]
\end{remark*}

\begin{remark*}[Predicate reading]
\[
\operatorname{PositiveFactorPreservesOrder}(\cdot_{\mathbb{Z}},\le_{\mathbb{Z}},0_{\mathbb{Z}},\mathbb{Z}).
\]
\end{remark*}

\begin{remark*}[Interpretation]
Multiplication by positive integers preserves non-strict order on both sides.
\end{remark*}

\begin{dependencies}
\begin{itemize}
  \item \hyperref[thm:left-positive-multiplication-preserves-order-on-z]{Left Multiplication by Positives Preserves Order on $\mathbb{Z}$}
  \item \hyperref[cor:right-positive-multiplication-preserves-order-on-z]{Right Multiplication by Positives Preserves Order on $\mathbb{Z}$}
\end{itemize}
\end{dependencies}

\begin{theorembox}{Theorem}
\begin{theorem}[Multiplication by Negatives Reverses Strict Order on $\mathbb{Z}$]
\label{thm:negative-multiplication-reverses-strict-order-on-z}
For all $x,y,z\in\mathbb{Z}$,
\[
z<0\land x<y \Longrightarrow z\cdot y<z\cdot x.
\]
\hyperref[prf:negative-multiplication-reverses-strict-order-on-z]{\textit{Go to proof.}}
\end{theorem}
\end{theorembox}

\begin{remark*}[Standard quantified statement]
\[
\forall x,y,z\in\mathbb{Z},\quad
z<0\land x<y \Longrightarrow z\cdot y<z\cdot x.
\]
\end{remark*}

\begin{remark*}[Predicate reading]
\[
\operatorname{NegativeFactorReversesOrder}(\cdot_{\mathbb{Z}},<_{\mathbb{Z}},0_{\mathbb{Z}},\mathbb{Z}).
\]
\end{remark*}

\begin{remark*}[Interpretation]
Multiplication by a negative integer reverses strict order.
\end{remark*}

\begin{dependencies}
\begin{itemize}
  \item \hyperref[thm:positive-multiplication-preserves-strict-order-on-z]{Multiplication by Positives Preserves Strict Order on $\mathbb{Z}$}
  \item \hyperref[cor:sign-rule-on-z]{Sign Rule on $\mathbb{Z}$}
  \item \hyperref[thm:strict-order-on-z-asymmetric]{Strict Order on $\mathbb{Z}$ Is Asymmetric}
\end{itemize}
\end{dependencies}

\begin{corollarybox}{Corollary}
\begin{corollary}[Multiplication by Negatives Reverses Order on $\mathbb{Z}$]
\label{cor:negative-multiplication-reverses-order-on-z}
For all $x,y,z\in\mathbb{Z}$,
\[
z<0\land x\le y \Longrightarrow z\cdot y\le z\cdot x.
\]
\hyperref[prf:negative-multiplication-reverses-order-on-z]{\textit{Go to proof.}}
\end{corollary}
\end{corollarybox}

\begin{remark*}[Standard quantified statement]
\[
\forall x,y,z\in\mathbb{Z},\quad
z<0\land x\le y \Longrightarrow z\cdot y\le z\cdot x.
\]
\end{remark*}

\begin{remark*}[Predicate reading]
\[
\operatorname{NegativeFactorReversesOrder}(\cdot_{\mathbb{Z}},\le_{\mathbb{Z}},0_{\mathbb{Z}},\mathbb{Z}).
\]
\end{remark*}

\begin{remark*}[Interpretation]
Multiplication by a negative integer reverses non-strict order.
\end{remark*}

\begin{dependencies}
\begin{itemize}
  \item \hyperref[thm:negative-multiplication-reverses-strict-order-on-z]{Multiplication by Negatives Reverses Strict Order on $\mathbb{Z}$}
  \item \hyperref[def:order-on-z]{Order on $\mathbb{Z}$}
\end{itemize}
\end{dependencies}

\begin{theorembox}{Theorem}
\begin{theorem}[Multiplication by Zero Collapses Order on $\mathbb{Z}$]
\label{thm:zero-multiplication-collapses-order-on-z}
For all $x,y\in\mathbb{Z}$,
\[
0\cdot x=0\cdot y.
\]
\hyperref[prf:zero-multiplication-collapses-order-on-z]{\textit{Go to proof.}}
\end{theorem}
\end{theorembox}

\begin{remark*}[Standard quantified statement]
\[
\forall x,y\in\mathbb{Z},\quad 0\cdot x=0\cdot y.
\]
\end{remark*}

\begin{remark*}[Predicate reading]
\[
\operatorname{ZeroMultiplicationCollapsesOrder}(\mathbb{Z},\cdot_{\mathbb{Z}},0_{\mathbb{Z}}).
\]
\end{remark*}

\begin{remark*}[Interpretation]
Multiplication by zero collapses all compared inputs to the same product.
\end{remark*}

\begin{dependencies}
\begin{itemize}
  \item \hyperref[cor:zero-absorption-on-z]{Zero Absorption on $\mathbb{Z}$}
  \item \hyperref[thm:multiplication-on-z-well-defined]{Multiplication on $\mathbb{Z}$ Is Well-Defined}
\end{itemize}
\end{dependencies}

\subsection*{\texorpdfstring{Absolute Value on $\mathbb{Z}$}{Absolute Value on Z}}

\begin{definitionbox}{Definition}
\begin{definition}[Absolute Value on $\mathbb{Z}$]
\label{def:absolute-value-on-z}
Define
\[
|x|=x\quad\text{if }0\le x,
\qquad
|x|=-x\quad\text{if }x<0.
\]
\end{definition}
\end{definitionbox}

\begin{remark*}[Standard quantified statement]
\[
(0\le x\Longrightarrow |x|=x)
\quad\text{and}\quad
(x<0\Longrightarrow |x|=-x).
\]
\end{remark*}

\begin{remark*}[Predicate reading]
\[
\operatorname{AbsoluteValue}(|\cdot|,\mathbb{Z},\le_{\mathbb{Z}},0_{\mathbb{Z}}).
\]
\end{remark*}

\begin{remark*}[Interpretation]
Absolute value chooses the nonnegative representative among an integer and its
additive inverse.
\end{remark*}

\begin{dependencies}
\begin{itemize}
  \item \hyperref[def:order-on-z]{Order on $\mathbb{Z}$}
  \item \hyperref[thm:trichotomy-on-z]{Trichotomy on $\mathbb{Z}$}
  \item \hyperref[thm:negation-on-z-well-defined]{Negation on $\mathbb{Z}$ Is Well-Defined}
\end{itemize}
\end{dependencies}

\begin{theorembox}{Theorem}
\begin{theorem}[Absolute Value Is Nonnegative]
\label{thm:absolute-value-nonnegative-on-z}
For every $x\in\mathbb{Z}$,
\[
0\le |x|.
\]
\hyperref[prf:absolute-value-nonnegative-on-z]{\textit{Go to proof.}}
\end{theorem}
\end{theorembox}

\begin{remark*}[Standard quantified statement]
\[
\forall x\in\mathbb{Z},\quad 0\le |x|.
\]
\end{remark*}

\begin{remark*}[Predicate reading]
\[
\operatorname{NonnegativeValue}(|\cdot|,\mathbb{Z},\le_{\mathbb{Z}},0_{\mathbb{Z}}).
\]
\end{remark*}

\begin{remark*}[Interpretation]
Absolute values are always nonnegative.
\end{remark*}

\begin{dependencies}
\begin{itemize}
  \item \hyperref[def:absolute-value-on-z]{Absolute Value on $\mathbb{Z}$}
  \item \hyperref[thm:trichotomy-on-z]{Trichotomy on $\mathbb{Z}$}
\end{itemize}
\end{dependencies}

\begin{theorembox}{Theorem}
\begin{theorem}[Absolute Value Vanishes Only at Zero]
\label{thm:absolute-value-zero-iff-on-z}
For every $x\in\mathbb{Z}$,
\[
|x|=0 \Longleftrightarrow x=0.
\]
\hyperref[prf:absolute-value-zero-iff-on-z]{\textit{Go to proof.}}
\end{theorem}
\end{theorembox}

\begin{remark*}[Standard quantified statement]
\[
\forall x\in\mathbb{Z},\quad |x|=0 \Longleftrightarrow x=0.
\]
\end{remark*}

\begin{remark*}[Predicate reading]
\[
\operatorname{ZeroKernel}(|\cdot|,\mathbb{Z},0_{\mathbb{Z}}).
\]
\end{remark*}

\begin{remark*}[Interpretation]
Only zero has absolute value zero.
\end{remark*}

\begin{dependencies}
\begin{itemize}
  \item \hyperref[def:absolute-value-on-z]{Absolute Value on $\mathbb{Z}$}
  \item \hyperref[thm:trichotomy-on-z]{Trichotomy on $\mathbb{Z}$}
\end{itemize}
\end{dependencies}

\begin{theorembox}{Theorem}
\begin{theorem}[Absolute Value Is Symmetric under Negation]
\label{thm:absolute-value-negation-on-z}
For every $x\in\mathbb{Z}$,
\[
|-x|=|x|.
\]
\hyperref[prf:absolute-value-negation-on-z]{\textit{Go to proof.}}
\end{theorem}
\end{theorembox}

\begin{remark*}[Standard quantified statement]
\[
\forall x\in\mathbb{Z},\quad |-x|=|x|.
\]
\end{remark*}

\begin{remark*}[Predicate reading]
\[
\operatorname{EvenUnderNegation}(|\cdot|,\mathbb{Z}).
\]
\end{remark*}

\begin{remark*}[Interpretation]
Absolute value ignores additive sign.
\end{remark*}

\begin{dependencies}
\begin{itemize}
  \item \hyperref[def:absolute-value-on-z]{Absolute Value on $\mathbb{Z}$}
  \item \hyperref[thm:trichotomy-on-z]{Trichotomy on $\mathbb{Z}$}
  \item \hyperref[cor:sign-rule-on-z]{Sign Rule on $\mathbb{Z}$}
\end{itemize}
\end{dependencies}

\begin{theorembox}{Theorem}
\begin{theorem}[Absolute Value Is Multiplicative]
\label{thm:absolute-value-multiplicative-on-z}
For all $x,y\in\mathbb{Z}$,
\[
|x\cdot y|=|x|\cdot |y|.
\]
\hyperref[prf:absolute-value-multiplicative-on-z]{\textit{Go to proof.}}
\end{theorem}
\end{theorembox}

\begin{remark*}[Standard quantified statement]
\[
\forall x,y\in\mathbb{Z},\quad |x\cdot y|=|x|\cdot |y|.
\]
\end{remark*}

\begin{remark*}[Predicate reading]
\[
\operatorname{Multiplicative}(|\cdot|,\mathbb{Z},\cdot_{\mathbb{Z}}).
\]
\end{remark*}

\begin{remark*}[Interpretation]
Absolute value turns products into products of absolute values.
\end{remark*}

\begin{dependencies}
\begin{itemize}
  \item \hyperref[def:absolute-value-on-z]{Absolute Value on $\mathbb{Z}$}
  \item \hyperref[prop:product-of-positives-is-positive]{Product of Positives Is Positive}
  \item \hyperref[cor:product-of-negatives-on-z]{Product of Negatives on $\mathbb{Z}$}
  \item \hyperref[thm:positive-multiplication-preserves-order-on-z]{Multiplication by Positives Preserves Order on $\mathbb{Z}$}
\end{itemize}
\end{dependencies}

\begin{theorembox}{Theorem}
\begin{theorem}[Bounding by Absolute Value]
\label{thm:absolute-value-bounds-on-z}
For every $x\in\mathbb{Z}$,
\[
-|x|\le x\le |x|.
\]
\hyperref[prf:absolute-value-bounds-on-z]{\textit{Go to proof.}}
\end{theorem}
\end{theorembox}

\begin{remark*}[Standard quantified statement]
\[
\forall x\in\mathbb{Z},\quad -|x|\le x\le |x|.
\]
\end{remark*}

\begin{remark*}[Predicate reading]
\[
\operatorname{BoundsByAbsoluteValue}(|\cdot|,\mathbb{Z},\le_{\mathbb{Z}}).
\]
\end{remark*}

\begin{remark*}[Interpretation]
An integer lies between the negative and positive forms of its absolute value.
\end{remark*}

\begin{dependencies}
\begin{itemize}
  \item \hyperref[def:absolute-value-on-z]{Absolute Value on $\mathbb{Z}$}
  \item \hyperref[thm:trichotomy-on-z]{Trichotomy on $\mathbb{Z}$}
  \item \hyperref[thm:absolute-value-negation-on-z]{Absolute Value Is Symmetric under Negation}
\end{itemize}
\end{dependencies}

\begin{theorembox}{Theorem}
\begin{theorem}[Triangle Inequality on $\mathbb{Z}$]
\label{thm:triangle-inequality-on-z}
For all $x,y\in\mathbb{Z}$,
\[
|x+y|\le |x|+|y|.
\]
\hyperref[prf:triangle-inequality-on-z]{\textit{Go to proof.}}
\end{theorem}
\end{theorembox}

\begin{remark*}[Standard quantified statement]
\[
\forall x,y\in\mathbb{Z},\quad |x+y|\le |x|+|y|.
\]
\end{remark*}

\begin{remark*}[Predicate reading]
\[
\operatorname{TriangleInequality}(|\cdot|,\mathbb{Z},+_{\mathbb{Z}},\le_{\mathbb{Z}}).
\]
\end{remark*}

\begin{remark*}[Interpretation]
Absolute value is subadditive for integer addition.
\end{remark*}

\begin{dependencies}
\begin{itemize}
  \item \hyperref[def:absolute-value-on-z]{Absolute Value on $\mathbb{Z}$}
  \item \hyperref[thm:absolute-value-bounds-on-z]{Bounding by Absolute Value}
  \item \hyperref[thm:addition-preserves-order-on-z]{Addition Preserves Order on $\mathbb{Z}$}
\end{itemize}
\end{dependencies}

\subsection*{\texorpdfstring{Well-Ordering Fails on $\mathbb{Z}$}{Well-Ordering Fails on Z}}

\begin{theorembox}{Theorem}
\begin{theorem}[$\mathbb{Z}$ Is Not Well-Ordered]
\label{thm:z-is-not-well-ordered}
The ordered set $(\mathbb{Z},\le)$ is not a well-order.
\hyperref[prf:z-is-not-well-ordered]{\textit{Go to proof.}}
\end{theorem}
\end{theorembox}

\begin{remark*}[Standard quantified statement]
\[
\neg\operatorname{WellOrder}(\le_{\mathbb{Z}},\mathbb{Z}).
\]
\end{remark*}

\begin{remark*}[Predicate reading]
\[
\operatorname{NotWellOrdered}(\mathbb{Z},\le_{\mathbb{Z}}).
\]
\end{remark*}

\begin{remark*}[Interpretation]
The integer order has infinite descending chains, so it is not a well-order.
\end{remark*}

\begin{dependencies}
\begin{itemize}
  \item \hyperref[def:well-order]{Well-Ordered Set}
  \item \hyperref[def:order-on-z]{Order on $\mathbb{Z}$}
  \item \hyperref[def:zero-in-z]{Zero in $\mathbb{Z}$}
  \item \hyperref[def:one-in-z]{One in $\mathbb{Z}$}
  \item \hyperref[thm:trichotomy-on-z]{Trichotomy on $\mathbb{Z}$}
\end{itemize}
\end{dependencies}
